\begin{proposition}[P4.14]\hypertarget{tag-02WL}{}\label{night:P4.14}\label{stmt:P4.14}
\StatementRecord{night:P4.14}{proved}{02WL}{the definitions and constructions stated in the proof}

The reducible-curve inference in \cite[proof of Proposition~7.3,
p.~16 of the attached-source edition]{CDP20} omits the two plane cubics
\[
 Z(Y^2-XZ)=0,\qquad XY(X-Y)=0.
\]
They are reduced Gorenstein curves of arithmetic genus one with
trivial dualizing sheaf and rational normalization, but have
respectively two tangent components and three concurrent
components. Their conductor order is two on every normalization
branch. These curves refute the asserted nodal-cycle
classification at that step; they are not counterexamples to
the threefold sign assertion.
\end{proposition}
\begin{proof}

\begin{enumerate}[label={$\langle1\rangle$\arabic*.},leftmargin=*]
\item Plane adjunction gives $\omega_C=\OO_C(3-3)=\OO_C$ and
$p_a(C)=1$. The first cubic has a line tangent to a smooth
conic at $[1:0:0]$. Its local equation is $y(y-x^2)=0$.
The normalization restrictions are $(p(t,0),p(t,t^2))$.
Their image is exactly the pairs $(a,b)$ with $b-a$ divisible
by $t^2$: the inverse expression is
$p(x,y)=a(x)+y(b(x)-a(x))/x^2$. Thus the normalization
quotient has length two and the conductor exponent is two on
each branch.

\item For the second cubic the branches are $(t,0),(0,t),(t,t)$.
All constant coefficients must agree, and the third linear
coefficient is the sum of the first two. These three conditions
are sufficient. In every degree $k\ge2$, the monomials
$x^k-x^{k-1}y$, $y^k-xy^{k-1}$ and $x^{k-1}y$ independently
realize the coefficients on the three branches. Hence the
normalization quotient has length three and each conductor
exponent is again two. The normalization Euler characteristics
are two and three; subtracting these quotient lengths gives
$\chi(\OO_C)=0$ in both cases.

\item The program \texttt{code/part4/genus\_one\_conductor.py}
executes the exact jet restriction matrices through orders
$2,3,4,6,8$, computing the ranks and conductor exponents.
When these branches vary in a surface, their permutation may
assemble them into a single normalization curve covering the
image conductor with degree two or three. Consequently a
unique horizontal conductor component with multiplicity two
does not alone force a degree-one map or generically reduced
image conductor. An additional argument is needed at that
precise point of the general irreducible case.

\end{enumerate}
\end{proof}